\section*{3.3 Gadget classification (author source pp.18--19)}
Now, we can prove a dichotomy result. First of all, we need the following classification of gadgets.
\begin{originalresult}{Lemma 3.8}
For every $g:\bits^m\to\bits$, one of the following cases holds:
\begin{enumerate}
\item[(i)] Function $g$ is a (possibly empty) conjunction or a disjunction of literals.
\item[(ii)] Function $g$ is weakly resistant.
\end{enumerate}
\end{originalresult}
\begin{proof}
Let $y_1,\ldots,y_m$ denote the inputs of $g$. If $g$ is not weakly resistant then there is a certificate $\alpha$ such that for every substitution $y_j=b$ conflicting with $\alpha$, the function $g|_{y_j=b}$ is constant. Moreover, we can assume that $\alpha$ is an inclusion-minimal certificate.
Suppose that $\alpha=\{y_{i_1}=b_1,\ldots,y_{i_t}=b_t\}$ and $g|_\alpha=b$. By assumption, the function $g|_{y_{i_s}=\bar b_s}$ is constant. Since $\alpha$ is inclusion-minimal, necessarily $g|_{y_{i_s}=\bar b_s}=\bar b$. Thus if $b=1$ then
\[g=[y_{i_1}=b_1]\wedge\cdots\wedge[y_{i_t}=b_t],\]
otherwise, if $b=0$ then
\[g=[y_{i_1}=b_1]\vee\cdots\vee[y_{i_t}=b_t].\]
\end{proof}
We can now state the dichotomy theorem.
\begin{originalresult}{Theorem 3.9}
For every gadget $g:\bits^m\to\bits$, one of the following cases holds:
\begin{enumerate}
\item[(i)] There is an infinite family of functions $f_n$ with $\DT(f_n)\to\infty$ and $\DTSize(f_n\comp g)=O(\DT(f_n))$.
\item[(ii)] For every function $f$, $\log\DTSize(f\comp g)=\Omega(\DT(f)^{\Omega(1)})$.
\end{enumerate}
\end{originalresult}
\begin{proof}
These two cases correspond to the cases of Lemma 3.8. If $g$ is a disjunction then we take $f_n:=\OR_n$. Clearly $\DT(f_n)=n$ and $\DTSize(f\comp g)=O(n)$ (if $g$ is constant, then $\DTSize(f_n)=1$). The case of a conjunction is similar.
Otherwise, $g$ is weakly resistant, thus by Corollary 3.5 we have $\log\DTSize(f\comp g)=\Omega(\sqrt{\DT(f)})$.
\end{proof}
